\documentclass[12pt,a4paper]{article}
\usepackage{amsmath,amssymb}
\usepackage{graphicx}
\usepackage{physics}
\usepackage{siunitx}
\usepackage{booktabs}
\usepackage{hyperref}
\usepackage[margin=1in]{geometry}

\title{Jeans Instability in IMBH-Cloud Tidal Disruption:\\
Tidal Heating vs. Radiative Cooling Balance}
\author{Analysis of Parameter Survey Simulations}
\date{\today}

\begin{document}
\maketitle

\begin{abstract}
We analyze the Jeans instability that caused the $r_p = 0.8$ pc isothermal cooling
simulation to crash. By examining the balance between tidal heating and radiative
cooling, we derive an estimate for the critical pericenter distance below which
tidal disruption dominates over gravitational collapse. We also calculate the
resolution requirements (particle number) needed to properly resolve the Jeans
mass at low temperatures.
\end{abstract}

\section{Introduction: The Problem}

The simulation with pericenter $r_p = 0.8$ pc crashed at $t \approx 2.88$ code units
due to runaway gravitational collapse. Analysis of the final snapshots reveals:
\begin{itemize}
    \item Maximum density: $\rho_{\max} \sim 10^7 \, M_\odot/\text{pc}^3$
          ($n \sim 2 \times 10^8 \, \text{cm}^{-3}$)
    \item Minimum temperature: $T_{\min} \sim 6$ K
    \item Jeans mass at collapse: $M_J \sim 0.06 \, M_\odot$
    \item Particle mass: $m_p \sim 0.008 \, M_\odot$
    \item Ratio: $M_J / m_p \sim 7$ (below resolution threshold)
\end{itemize}

\section{Theoretical Framework}

\subsection{Jeans Mass and Length}

The Jeans mass is the critical mass above which a cloud will collapse under
self-gravity:
\begin{equation}
    M_J = \frac{\pi^{5/2}}{6} \frac{c_s^3}{G^{3/2} \rho^{1/2}}
\end{equation}
where $c_s = \sqrt{k_B T / (\mu m_H)}$ is the isothermal sound speed.

For molecular gas with $\mu = 1.27$:
\begin{equation}
    \boxed{M_J \approx 18 \, M_\odot \left(\frac{T}{10 \, \text{K}}\right)^{3/2}
    \left(\frac{n}{10^4 \, \text{cm}^{-3}}\right)^{-1/2}}
\end{equation}

The Jeans length is:
\begin{equation}
    \lambda_J = \sqrt{\frac{\pi c_s^2}{G \rho}} \approx 0.4 \, \text{pc}
    \left(\frac{T}{10 \, \text{K}}\right)^{1/2}
    \left(\frac{n}{10^4 \, \text{cm}^{-3}}\right)^{-1/2}
\end{equation}

\subsection{Tidal Heating Rate}

As the cloud approaches the IMBH, tidal forces do work on the gas. The tidal
heating rate per unit mass can be estimated as:
\begin{equation}
    \dot{\epsilon}_{\text{tidal}} \sim \frac{G M_{\text{BH}}}{r^3} R_{\text{cloud}}^2
    \cdot v_{\text{rel}}
\end{equation}
where $v_{\text{rel}}$ is the relative velocity and $r$ is the distance to the IMBH.

At pericenter passage ($r = r_p$, $v = v_p$):
\begin{equation}
    \dot{\epsilon}_{\text{tidal}} \sim \frac{G M_{\text{BH}} R^2 v_p}{r_p^3}
\end{equation}

For a parabolic orbit, $v_p = \sqrt{2 G M_{\text{BH}} / r_p}$, so:
\begin{equation}
    \boxed{\dot{\epsilon}_{\text{tidal}} \sim \frac{\sqrt{2} G^{3/2} M_{\text{BH}}^{3/2} R^2}{r_p^{7/2}}}
\end{equation}

\textbf{Key scaling:} $\dot{\epsilon}_{\text{tidal}} \propto r_p^{-7/2}$

\subsection{Radiative Cooling Rate}

The Koyama \& Inutsuka (2000) cooling function gives:
\begin{equation}
    \Lambda(T) = 2 \times 10^{-19} \exp\left(-\frac{1.184 \times 10^5}{T + 1000}\right)
    + 2.8 \times 10^{-28} \sqrt{T} \exp\left(-\frac{92}{T}\right)
    \, \text{erg cm}^3 \text{s}^{-1}
\end{equation}

The cooling rate per unit mass is:
\begin{equation}
    \dot{\epsilon}_{\text{cool}} = \frac{n^2 \Lambda(T)}{\rho} = \frac{n \Lambda(T)}{\mu m_H}
\end{equation}

For $T \sim 35$ K and $n \sim 500 \, \text{cm}^{-3}$:
\begin{equation}
    \dot{\epsilon}_{\text{cool}} \sim 10^{-3} \, \text{erg g}^{-1} \text{s}^{-1}
\end{equation}

\section{Balance Condition: Critical Pericenter}

\subsection{Heating-Cooling Equilibrium}

The cloud remains thermally stable when tidal heating balances cooling:
\begin{equation}
    \dot{\epsilon}_{\text{tidal}} \gtrsim \dot{\epsilon}_{\text{cool}}
\end{equation}

Substituting our expressions:
\begin{equation}
    \frac{G^{3/2} M_{\text{BH}}^{3/2} R^2}{r_p^{7/2}} \gtrsim \frac{n \Lambda(T)}{\mu m_H}
\end{equation}

Solving for critical pericenter:
\begin{equation}
    \boxed{r_{p,\text{crit}} \lesssim \left(\frac{G^{3/2} M_{\text{BH}}^{3/2} R^2 \mu m_H}{n \Lambda(T)}\right)^{2/7}}
\end{equation}

\subsection{Numerical Estimate}

Using simulation parameters:
\begin{align}
    M_{\text{BH}} &= 10^5 \, M_\odot \\
    R_{\text{cloud}} &= 1.55 \, \text{pc} \\
    n &= 500 \, \text{cm}^{-3} \\
    T &= 35 \, \text{K} \\
    \Lambda(35 \, \text{K}) &\approx 10^{-27} \, \text{erg cm}^3 \text{s}^{-1}
\end{align}

Converting to code units and computing:
\begin{equation}
    r_{p,\text{crit}} \approx 0.5 - 0.6 \, \text{pc}
\end{equation}

\subsection{Interpretation}

\begin{center}
\begin{tabular}{lll}
\toprule
Pericenter & Regime & Outcome \\
\midrule
$r_p < r_{p,\text{crit}}$ & Tidal heating dominates & Cloud disrupted, stable \\
$r_p > r_{p,\text{crit}}$ & Cooling dominates & Jeans collapse, unstable \\
\midrule
$r_p = 0.4$ pc & $< r_{p,\text{crit}}$ & \textbf{Stable} (observed) \\
$r_p = 0.8$ pc & $> r_{p,\text{crit}}$ & \textbf{Collapsed} (observed) \\
\bottomrule
\end{tabular}
\end{center}

This explains why:
\begin{itemize}
    \item $r_p = 0.4$ pc simulation with cooling remains stable
    \item $r_p = 0.8$ pc simulation with cooling collapsed
\end{itemize}

\section{Physical Reality: Jeans Floor and Opacity}

\subsection{Why Real Clouds Don't Collapse to Singularities}

In reality, several physical mechanisms prevent infinite collapse:

\subsubsection{1. Opacity-Limited Cooling}

At high densities ($n \gtrsim 10^{10}$ cm$^{-3}$), the gas becomes optically thick
to its own cooling radiation. The cooling timescale becomes:
\begin{equation}
    t_{\text{cool}} \sim \frac{\tau}{n \Lambda} \propto \rho
\end{equation}
where $\tau$ is the optical depth. This effectively reduces cooling efficiency
at high densities.

\subsubsection{2. Dust Thermal Coupling}

At $n \gtrsim 10^4$ cm$^{-3}$, gas-dust collisions thermally couple the gas to
dust grains. The dust temperature is set by the interstellar radiation field:
\begin{equation}
    T_{\text{dust}} \sim 10 \, \text{K} \cdot \left(\frac{G_0}{1}\right)^{1/6}
\end{equation}
This provides a temperature floor $T \gtrsim 10$ K.

\subsubsection{3. Compressional Heating}

Rapid collapse ($t_{\text{collapse}} < t_{\text{cool}}$) leads to adiabatic heating:
\begin{equation}
    T \propto \rho^{\gamma - 1}
\end{equation}
When collapse becomes faster than cooling, the gas heats up, increasing pressure
support and slowing collapse.

\subsection{Effective Equation of State}

The combination of these effects creates an effective ``stiff'' equation of state
at high densities:
\begin{equation}
    P = K \rho^{\gamma_{\text{eff}}}
\end{equation}
where $\gamma_{\text{eff}} > 4/3$ above a critical density, preventing runaway
collapse.

\subsection{Implementing a Jeans Floor}

The Jeans floor is a numerical technique that mimics these physical effects by
preventing cooling below a temperature where the Jeans mass is resolved:
\begin{equation}
    T_{\text{floor}}(\rho) = T_J \left(\frac{\rho}{\rho_J}\right)^{2/3}
\end{equation}
where $T_J$ and $\rho_J$ are chosen such that $M_J(T_{\text{floor}}, \rho) =
N_{\text{neighbor}} \cdot m_{\text{particle}}$.

This ensures:
\begin{equation}
    \boxed{M_J \geq N_{\text{neighbor}} \cdot m_p \sim 50 \cdot m_p}
\end{equation}

\section{Resolution Requirements}

\subsection{Jeans Criterion for SPH}

For SPH to properly capture gravitational collapse without artificial
fragmentation, the Jeans mass must be resolved by at least $N_J$ particles:
\begin{equation}
    M_J \geq N_J \cdot m_p
\end{equation}

The standard criterion is $N_J \approx 2 N_{\text{neighbor}} \approx 100$
(Bate \& Burkert 1997).

\subsection{Required Particle Number}

Given the observed conditions at collapse:
\begin{align}
    T_{\min} &= 6 \, \text{K} \\
    n_{\max} &= 2 \times 10^8 \, \text{cm}^{-3} \\
    M_J &= 18 \cdot (6/10)^{1.5} \cdot (2 \times 10^8 / 10^4)^{-0.5} \\
        &= 18 \cdot 0.465 \cdot 0.00707 \\
        &\approx 0.059 \, M_\odot
\end{align}

For the cloud mass $M_{\text{cloud}} = 127.5 \, M_\odot$:

\textbf{Current resolution:}
\begin{equation}
    N_{\text{current}} = 15552 \implies m_p = 0.0082 \, M_\odot
\end{equation}
\begin{equation}
    \frac{M_J}{m_p} = \frac{0.059}{0.0082} \approx 7.2 \quad \text{(INSUFFICIENT)}
\end{equation}

\textbf{Required resolution for $N_J = 100$:}
\begin{equation}
    m_p^{\text{req}} = \frac{M_J}{N_J} = \frac{0.059}{100} = 5.9 \times 10^{-4} \, M_\odot
\end{equation}
\begin{equation}
    \boxed{N_{\text{required}} = \frac{M_{\text{cloud}}}{m_p^{\text{req}}}
    = \frac{127.5}{5.9 \times 10^{-4}} \approx 2.2 \times 10^5}
\end{equation}

\subsection{Scaling with Temperature}

If we use a Jeans floor at $T_{\text{floor}} = 10$ K:
\begin{equation}
    M_J(10 \, \text{K}, 10^8 \, \text{cm}^{-3}) = 18 \cdot 1 \cdot 0.01 = 0.18 \, M_\odot
\end{equation}
\begin{equation}
    N_{\text{required}} = \frac{127.5 \cdot 100}{0.18} \approx 7 \times 10^4
\end{equation}

With $T_{\text{floor}} = 20$ K:
\begin{equation}
    M_J(20 \, \text{K}, 10^8 \, \text{cm}^{-3}) = 18 \cdot 2.83 \cdot 0.01 = 0.51 \, M_\odot
\end{equation}
\begin{equation}
    N_{\text{required}} = \frac{127.5 \cdot 100}{0.51} \approx 2.5 \times 10^4
\end{equation}

\subsection{Summary Table}

\begin{center}
\begin{tabular}{cccc}
\toprule
$T_{\text{floor}}$ (K) & $M_J$ ($M_\odot$) & $N_{\text{required}}$ & Feasible? \\
\midrule
6 (no floor) & 0.059 & $2.2 \times 10^5$ & Expensive \\
10 & 0.18 & $7 \times 10^4$ & Moderate \\
20 & 0.51 & $2.5 \times 10^4$ & Current resolution \\
35 (initial) & 1.4 & $9 \times 10^3$ & Easy \\
\bottomrule
\end{tabular}
\end{center}

\section{Conclusions}

\begin{enumerate}
    \item The $r_p = 0.8$ pc simulation crashed due to Jeans instability when
          cooling reduced $T$ below the Jeans resolution limit.

    \item The critical pericenter is $r_{p,\text{crit}} \approx 0.5-0.6$ pc,
          explaining why $r_p = 0.4$ pc is stable and $r_p = 0.8$ pc is not.

    \item Physical mechanisms (opacity, dust coupling, adiabatic heating)
          naturally provide temperature floors in real ISM.

    \item To resolve Jeans collapse at $T = 6$ K without a floor, we need
          $N \sim 2 \times 10^5$ particles.

    \item A practical solution is implementing a Jeans floor at $T \sim 10-20$ K,
          which requires only $N \sim 2.5-7 \times 10^4$ particles.
\end{enumerate}

\section{Recommended Action}

Implement a density-dependent temperature floor:
\begin{equation}
    T_{\text{floor}}(\rho) = \max\left(10 \, \text{K}, \, T_J \left(\frac{\rho}{\rho_J}\right)^{2/3}\right)
\end{equation}
where $T_J = 10$ K and $\rho_J$ is chosen such that $M_J(T_J, \rho_J) = 50 \cdot m_p$.

This is physically motivated and computationally efficient.

\end{document}
